\begin{abstract}
We ask whether a conditional diffusion policy over action chunks controls a simulated
articulated robot more reliably than deterministic and recurrent behavior cloning when all
of them learn from the same demonstrations. On the robosuite Lift and PickPlaceCan tasks
(MuJoCo physics, a Panda arm with operational-space control) we train four policy families
on the robomimic proficient-human demonstrations with exact simulator state as input: a
one-step MLP, an LSTM, an MLP that predicts action chunks, and a 1D temporal U-Net
diffusion policy, with matched parameter counts (\rf{params.bcmlp} to \rf{params.chunk}),
\rf{proto.nseeds} training seeds
and a matched, preregistered tuning budget. Checkpoints are chosen on validation rollouts,
and the frozen protocol is then run once on \rf{proto.nstates} held-out initial states per
task. On Can, diffusion succeeds in \rf{l1.can.diff.succ}\pct{} of held-out episodes against
\rf{l1.can.bcmlp.succ}\pct{} for the one-step MLP (paired difference
\rf{l1.can.diff-bcmlp.d}\pp, \rf{proto.level}\pct{} interval
[\rf{l1.can.diff-bcmlp.lo}, \rf{l1.can.diff-bcmlp.hi}]), although a seed-level $t$
interval, computed after the audit and not preregistered, includes zero; on Lift the same
difference is not separated by the preregistered bootstrap (the seed-level interval does
separate it). Diffusion is not separated from the recurrent policy or from the
chunked MLP on either task, so these data give no evidence that the generative model adds to
observation history and action chunking, while diffusion costs
\rf{l1.ratio.cpu.lo}--\rf{l1.ratio.cpu.hi}$\times$ the chunked MLP's training time.
Under preregistered perturbations of the initial arm configuration, diffusion degrades least
on Lift, but on Can its drop is not separated from the chunked MLP's
(\rf{rob.can.dd.arm.diff-chunk.d}\pp{} [\rf{rob.can.dd.arm.diff-chunk.lo},
\rf{rob.can.dd.arm.diff-chunk.hi}]); the diffusion sampler's bound on its actions is a
confound in these comparisons. Two separately frozen extensions follow: on Lift, visual diffusion and visual
chunked BC from two cameras both reach the task's ceiling, and a DP3-inspired point-cloud
diffusion policy generalizes worse than the frozen RGB diffusion policy (same head and
training, different encoder and augmentation) to cube placements
outside the demonstrations' range (\rf{l2.h5.d}\pp{} [\rf{l2.h5.lo}, \rf{l2.h5.hi}]). The
main limitation is resolution: \rf{proto.nseeds} seeds and \rf{proto.nstates} starts per
task cannot resolve differences of a few points near the ceiling, and everything here is
simulation.
\end{abstract}
